\documentclass[10pt, letterpaper]{article}

\usepackage{../../../../template/template}

\setDocTitle{Verbale interno 09/03/2026}
\setDocVersion{0.1.0}
\setDocData{09/03/2026}
\setDocIndex{vi\_2026\_03\_09}
\setDocState{1}
\setDocRecipient{1}

\begin{document}

\makeTitlePage

% \newpage
% \addChangelog{0.1.0}{11/03/2026}{F. Pasqual}{C. Libralato}{\noOne}{\firstDraft}
% \makeChangelog

\newpage

\tableofcontents

\newpage

% -------------------------------------------------
% Informazioni riunione
% -------------------------------------------------
\makeMeetingInfoTable{09/03/2026}{14:30}{16:30}{via \textit{Discord$_G$}}

\addParticipant{Valeria}{Baleanu}{Progettista}{P}
\addParticipant{Leonardo}{Pellizzon}{Progettista}{P}
\addParticipant{Luca}{Granziero}{Progettista}{P}
\addParticipant{Francesco}{Pasqual}{Responsabile}{P}
\addParticipant{Giuseppe}{De Fina}{Progettista}{P}
\addParticipant{Christian}{Libralato}{Verificatore}{P}
\addParticipant{Filippo}{Venzo}{Amministratore}{P}
\makeMeetingParticipantsTable


\addPrecedentTodo{-}{\noOne}{Manca il verbale interno della scorsa riunione, aspetto Luca}{-}


\makePrecedentTodoTable
% -------------------------------------------------
% Ordine del giorno
% -------------------------------------------------
\section{Ordine del giorno}
\begin{itemize}
    \item Coordinamento ruoli progettista/programmatore e aggiornamento \textit{sprint$_G$};
    \item gestione e design \textit{backend}: allarmi, notifiche e gestione utenti;
    \item organizzazione \textit{frontend} e personalizzazione dashboard;
    \item gestione documentazione e strumenti di supporto;
    \item divisione compiti e pianificazione ultime attività.
\end{itemize}

% -------------------------------------------------
% Approfondimento
% -------------------------------------------------
\section{Approfondimento}

\subsection*{Coordinamento ruoli e aggiornamento \textit{sprint}}
Il team ha discusso la necessità di completare interamente la fase di progettazione prima di passare alla programmazione effettiva.
Per garantire questo risultato, è stato deciso che durante la prima metà del prossimo sprint anche i membri con ruoli diversi dal progettista lavoreranno in veste di progettisti nei momenti di necessità.
Si è inoltre stabilito di mantenere inalterata la durata degli \textit{sprint} a due settimane per minimizzare le perdite di tempo legate all'overhead gestionale.
Finalmente, è emersa la necessità di aggiornare il Piano di Progetto e il Piano di Qualifica includendo dati, metriche e grafici aggiornati fino allo \textit{sprint} 9.

\subsection*{Gestione e design backend}
Sono stati esaminati i progressi lato \textit{backend}, focalizzandosi in particolare sulla modellazione dettagliata
degli allarmi e dei relativi \textit{use case$_G$}, oltre che sulla gestione degli utenti e sulla sincronizzazione \textit{frontend}/\textit{backend} tramite \textit{DTO$_G$}.
Un punto cruciale della discussione ha riguardato l'importanza di differenziare nel \textit{database} le chiavi naturali da quelle surrogate per le entità di reparti,
utenti e dispositivi, tema che richiederà un ulteriore approfondimento con i docenti e il referente aziendale.
Sono state inoltre affrontate l'implementazione della comunicazione via \textit{WebSocket$_G$} per le notifiche rapide, l'applicazione rigorosa dei principi \textit{SOLID$_G$} e la messa in sicurezza degli \textit{endpoint$_G$} con logiche specifiche per i \textit{token JWT$_G$}.

\subsection*{Organizzazione frontend e dashboard}
Per quanto concerne il \textit{frontend}, la discussione si è concentrata sulla struttura generale e sulla necessità
di definire con precisione i servizi a livello 3 del modello \textit{C4}.
È stata avanzata l'ipotesi di rendere la \textit{dashboard$_G$} personalizzabile, lasciando come punto in sospeso se memorizzare tali preferenze lato utente o lato amministratore.
A supporto della progettazione grafica e del design collaborativo è stato confermato l'impiego di \textit{Figma$_G$}.
Resta da verificare la disponibilità e l'applicabilità di eventuali pacchetti \textit{Angular} integrabili per il \textit{wireframing} e la generazione di codice.

\subsection*{Gestione documentazione e strumenti di supporto}
Analizzando le limitazioni dell'attuale strumento utilizzato per i diagrammi e l'esigenza di un miglior
versionamento, è stato deciso di archiviare i file dei diagrammi e i documenti di specifica direttamente su \textit{GitHub$_G$},
centralizzando così le risorse e facilitandone l'accessibilità.

\subsection*{Divisione compiti e pianificazione ultime attività}
La riunione si è conclusa con una chiara ripartizione dei carichi di lavoro per le fasi imminenti:
considerate le imminenti indisponibilità esterne di alcuni componenti, è stata ribadita la priorità assoluta di implementare prima le funzionalità obbligatorie del progetto, posticipando quelle opzionali.
È stato confermato l'impegno a mantenere una comunicazione costante e a sfruttare appieno gli strumenti collaborativi per portare a termine i lavori nei tempi previsti.
Infine, si è ricordata la necessità di redarre le presentazioni per il Diario di Bordo e per il \textit{SAL$_G$} con l'azienda del 11/03/2026.

% -------------------------------------------------
% Decisioni
% -------------------------------------------------
\addDecision{Dedicare la prima metà dello sprint esclusivamente alla progettazione con tutti i membri operativi come progettisti}
\addDecision{Aggiornare i dati e le metriche del Piano di Progetto e Piano di Qualifica fino allo sprint 9}
\addDecision{Mantenere gli sprint della durata di due settimane}
\addDecision{Utilizzo di Figma per il design grafico collaborativo}
\addDecision{Centralizzare la documentazione e i diagrammi SVG dei diagrammi delle classi su GitHub per migliorarne versionamento e gestione}
\makeDecisionTable

% -------------------------------------------------
% Azioni
% -------------------------------------------------
\addTodo{\#161}{Stesura verbale interno relativo alla riunione del 09/03/2026}{F. Pasqual, C. Libralato}{16/03/2026}
\addTodo{\#163}{Redazione presentazione per Diario di Bordo}{F. Pasqual, C. Libralato}{11/03/2026}
\addTodo{\#164}{Aggiornare Piano di Progetto fino allo sprint 9}{F. Venzo}{22/03/2026}
\addTodo{\#166}{Redazione presentazione per \textit{SAL}}{F. Pasqual, C. Libralato}{11/03/2026}
\addTodo{\#167}{Preparare bozzetti grafici e design frontend collaborativo su Figma}{V. Baleanu, L. Pellizzon, L. Granziero, G. De Fina}{18/03/2026}
\addTodo{\#165}{Aggiornare Piano di Qualifica fino allo sprint 9}{F. Venzo}{22/03/2026}
\addTodo{\#169}{Aggiungere nuovi termini al glossario}{F. Pasqual}{16/03/2026}
\addTodo{\#172}{Completare modellazione classi backend di notifiche, WebSocket, allarmi e utenti}{V. Baleanu, L. Pellizzon, L. Granziero, G. De Fina}{16/03/2026}
\addTodo{\#174}{Completare modellazione classi frontend}{V. Baleanu, L. Pellizzon, L. Granziero, G. De Fina}{16/03/2026}
\addTodo{\#173}{Creare documentazione condivisa su GitHub con caricamento diagrammi e specifiche}{F. Pasqual}{16/03/2026}

\makeTodoTable

\end{document}